\section{Proofs of main results}
\label{main_results_section}


































\subsection{Main tool} \label{reg_ss}

We begin with our main tool, Theorem~\ref{main_lemma_in_intro}, a result that allows us to evaluate primary characters on regular semisimple elements more easily.
We present some lemmas before giving the proof.
In the preliminary lemmas and in Theorem~\ref{main_lemma_in_intro}, the hypotheses include ``\(\mu \vdash n\) and \(g \in \CT_\mu^\Box (q)\).''
In all the proofs, we will assume \(g\) is in rational canonical form and denote the distinct irreducible factors of the characteristic polynomial of \(g\) by \(h_1,\ldots,h_{\ell(\mu)}\) with \(\deg h_i = \mu_i\) for each \(i \in \{1,\ldots,\ell(\mu)\}\).
Note that this implies each diagonal block \(\pi_i^\mu (g) \in \GL_{\mu_i} (q)\) is in the primary conjugacy class indexed by \(h_i \mapsto \Box\).


\begin{lemma}
\label{d_divides_mu}
For all \(n \in \N\), \(d \mid n\), \(\nu \vdash n / d\), \(\mu \vdash n\), prime powers \(q\), \(g \in \CT_\mu^\Box (q)\), and \(\alpha : d \Z \to \Z\), we have \(B_{d\nu}^\alpha(g) = 0\) unless \(d \nu = \mu\).
\end{lemma}
\begin{proof}
By definition,
\begin{equation}
\label{eqn_for_d_divides_mu}
B_{d \nu}^\alpha (g)
=
\frac{1}{\#\para_{d \nu}}
\sum_{\substack{x \in \GL_n \F_q \\ xgx^{-1} \in \para_{d \nu}}}
\prod_{i = 1}^{\ell(\nu)}
P_{d \nu_i}^{\alpha (d \nu_i)} \left (\pi^{d \nu}_i (xgx^{-1}) \right ).
\end{equation}
Consider a single summand in \eqref{eqn_for_d_divides_mu}, which is a product of Primary-support character values.
By definition, the Primary-support characters vanish away from primary conjugacy classes.
Thus, the product of them vanishes if any diagonal block of \(xgx^{-1}\) is not primary.
So assume the product of the $P$ characters in \eqref{eqn_for_d_divides_mu} does not vanish, and hence each block \(\pi_i^{d \nu} (xgx^{-1})\) is primary.

For each \(i \in \{1,\ldots,\ell(\nu)\}\), let \(\tilde{h}_i\) be the characteristic polynomial of \(\pi^{d \nu}_i (xgx^{-1})\).
Since \(xgx^{-1} \in \para_{d \nu}\) has a block-upper-triangular structure, its characteristic polynomial equals \(\prod_{i=1}^{\ell(\nu)} \tilde{h}_i\).
The fact that each \(\pi^{d \nu}_i ( xgx^{-1})\) is primary implies that each \(\tilde{h}_i\) is a power \(\rho_i^{a_i}\) of an irreducible polynomial \(\rho_i \in \CF (q)\).
On the other hand, \(g\) is regular semisimple, meaning its characteristic polynomial has no repeated factors.
Thus, each \(a_i = 1\) and there exists a permutation \(\sigma \in \FS_{\ell(\mu)}\) such that \(\rho_i = \tilde{h}_i = h_{\sigma(i)}\) for all \(i \in \{1,\ldots,\ell(\mu)\}\).
Computing degrees show that \(d \nu_i = \deg \tilde{h}_i = \deg h_{\sigma(i)} = \mu_{\sigma(i)}\) for all \(i \in \{1,\ldots,\ell(\mu)\}\).
This implies \(d \nu = \mu\).
\end{proof}



We require some more terminology before moving forward.
Given \(\mu \vdash n\), refer to a flag $S_\bullet$ of nested subspaces
$$ S_1 \subset S_2 \subset \cdots \subset S_{\ell(\mu)}$$
of $V$
as a \myemph{$\mu$-flag} if
$$\dim S_j = \sum_{i=1}^j \mu_i$$ for all $j \in \{1,\ldots,\ell(\mu)\}$.
Refer to an ordered basis $(e_1,\ldots,e_{n})$ of $V$ as a \myemph{basis for $S_\bullet$} if
$$ \left (e_1,\ldots,e_{\sum_{i=1}^j \mu_i} \right ) $$
is a basis for $S_j$ for each $j \in \{1,\ldots,\ell(\mu)\}$.
Conversely, each ordered basis \((e_1,\ldots,e_n)\) for \(V\) determines a \(\mu\)-flag by taking the \(j^\text{th}\) subspace in the flag to be the span of \(\left ( e_1,\ldots,e_{\sum_{i=1}^j \mu_i} \right )\) for each \(j \in \{1,\ldots,\ell(\mu)\}\).
Given a $\mu$-flag $S_\bullet$, say that a matrix in $\GL_{n} \F_q$ \myemph{stabilizes $S_\bullet$} if it stabilizes $S_j$ for each $j \in \{1,\ldots,\ell(\mu)\}$.

\begin{lemma}
\label{parabolic_counting_lemma}
For all \(n \in \N\), $\mu \vdash n$, prime powers \(q\), and $g \in \CT_\mu^\Box (q)$, we have
\begin{equation} \label{parabolic_counting_eqn} \# \left \{ x \in \GL_n \F_q : xgx^{-1} \in \para_\mu \right \} = \# \para_\mu \cdot \prod_{i \geq 1} m_i(\mu)! .\end{equation}
\end{lemma}
\begin{proof}
Viewing \(g\) abstractly as a linear transformation on \(V\), the left side of \eqref{parabolic_counting_eqn} is the number of ordered bases of $V$ with respect to which the matrix representing $g$ is an element of $\para_\mu$.
For any fixed basis $\CB = (v_1,\ldots,v_{n})$ for $V$, being an element of $\para_\mu$ is equivalent to stabilizing the $\mu$-flag determined by $\CB$.
Therefore, the left side of \eqref{parabolic_counting_eqn} is the product of
\begin{enumerate}[(1)]
\item the number of $\mu$-flags that $g$ stabilizes and
\item the number of ordered bases for a given $\mu$-flag.
\end{enumerate}
The second part is $\# \para_\mu$.

For the first part, consider a $\mu$-flag
$$ S_\bullet = S_1 \subset S_2 \subset \cdots \subset S_{\ell(\mu)}$$
that $g$ stabilizes.
For each \(i \in \{1,\ldots,\ell(\mu)\}\), let \(Q_i\) denote the quotient \(S_i / S_{i-1}\), with the convention that \(S_0\) is zero-dimensional.
Note that \(\dim Q_i = \mu_i\).
Under the isomorphism \eqref{rs_rcf}, let \(E_i \subset V\) denote the subspace corresponding to \(\F_q[z] / (h_i(z))\) for each \(i \in \{1,\ldots,\ell(\mu)\}\).
By Corollary \ref{really_helpful}, the only subspaces of $V$ that $g$ stabilizes are direct sums of the $E_i$'s.
Therefore, each $S_j$ is a direct sum of the $E_i$'s.
This implies each $Q_i$ is also a direct sum of the $E_i$'s.
Working backwards from \(Q_{\ell(\mu)}\) to \(Q_1\), we see that
\(
Q_{\ell(\mu)-m_1(\mu)+1},
\ldots,
Q_{\ell(\mu)}
\) are all 1-dimensional
and thus form a permutation of the \(m_1(\mu)\) subspaces
\(
E_{\ell(\mu)-m_1(\mu)+1},
\ldots,
E_{\ell(\mu)}
\).
Likewise,
\[
Q_{\ell(\mu)-m_1(\mu)-m_2(\mu)+1},
\ldots,
Q_{\ell(\mu)-m_1(\mu)}
\]
are all 2-dimensional and thus form a permutation of the \(m_2(\mu)\) subspaces
\[
E_{\ell(\mu)-m_1(\mu)-m_2(\mu)+1},
\ldots,
E_{\ell(\mu)-m_1(\mu)},
\]
as there are no 1-dimensional \(E_i\)'s remaining.
Continuing in this fashion, we see that the sequence \((Q_1,\ldots,Q_{\ell(\mu)})\) of quotients is one of \(\prod_{i \geq 1} m_i(\mu)!\) total possibilities.
Since the quotients determine \(S_\bullet\) uniquely, the result follows.
\end{proof}

We are now ready to prove Theorem~\ref{main_lemma_in_intro}.
For \(d \mid n\), \(f \in \CF_d (q)\), and \(\lambda \vdash n/d\), it states that,
if some part of \(\mu\) is not divisible by \(d\), then \(\chi^{f \mapsto \lambda} (g) = 0\), and otherwise, there exists \(\tilde{\mu} \vdash n/d\) such that \(\mu = d \tilde{\mu}\), and we have
\begin{equation}
\chi^{f \mapsto \lambda} (g)
= (-1)^{\tfrac{n}{d}(d-1)}
\chi^\lambda_{\tilde{\mu}}
\prod_{i = 1}^{\ell(\mu)}
\frac{1}{\tilde{\mu}_i}
\sum_{\substack{ \beta_i \in \F_{q^{\mu_i}} \\ h_i(\beta_i) = 0}} \te (\beta_i)^{\ell_f \cdot [\tilde{\mu}_i]_{q^d}}.
\end{equation}



\begin{proof}[Proof of Theorem \ref{main_lemma_in_intro}]
By definitions \eqref{jrreducible} and \eqref{grand_finale}, we have
\begin{equation}
\label{in_the_beginning}
\chi^{f \mapsto \lambda} (g)
= J_f^\lambda (g) = (-1)^{\tfrac{n}{d} (d-1)} \cdot \sum_{\nu \vdash \tfrac{n}{d}} \frac{\chi^\lambda_\nu}{z_\nu} B_{d \nu}^{\ell_f \cdot \alpha_{d}} (g).
\end{equation}
By Lemma \ref{d_divides_mu}, $B_{d \nu}^{\ell_f \cdot \alpha_d} (g) = 0$ unless $d \nu = \mu$.
If some part of $\mu$ is not divisible by $d$, then $\chi^{f \mapsto \lambda}(g) = 0$, proving the first statement in the lemma.
Otherwise, there exists a unique partition \(\tilde{\mu} \vdash n/d\) such that \(\mu = d \tilde{\mu}\),
and only the summand corresponding to \(\tilde{\mu}\) in \eqref{in_the_beginning} does not vanish.
In this case, \eqref{in_the_beginning} reduces to
\begin{equation}
\label{super_concise}
\chi^{f \mapsto \lambda} (g) = (-1)^{\frac{n}{d}(d-1)} \frac{\chi^\lambda_{\tilde{\mu}}}{z_{\tilde{\mu}}} B_{\mu}^{\ell_f \cdot \alpha_d} (g).
\end{equation}
By definitions \eqref{parabolic_induction} and \eqref{B_chars}, we can rewrite \eqref{super_concise} as
\begin{equation}
\label{expandedB_eqn}
\chi^{f \mapsto \lambda} (g)
= (-1)^{\tfrac{n}{d} (d-1)} \cdot \frac{\chi^\lambda_{\tilde{\mu}}}{z_{\tilde{\mu}}} \cdot \frac{1}{\# \para_{\mu}} \sum_{\substack{x \in \GL_n \F_q \\ xgx^{-1} \in \para_{\mu}}} \prod_{i = 1}^{\ell(\mu)} P_{\mu_i}^{\ell_f \cdot \alpha_d (\mu_i)} (\pi^\mu_i (xgx^{-1})).
\end{equation}

Consider the summation in \eqref{expandedB_eqn}.
It can be rewritten as
\begin{equation}
\label{expandedB_eqn_summation}
\sum_{\substack{x \in \GL_n \F_q \\ xgx^{-1} \in \para_\mu}}
\prod_{s \geq 1}
\prod_{\{j \in \N \, : \, \mu_j = s\}}
P_s^{\ell_f \cdot \alpha_d(s)}
(\pi_j^\mu (xgx^{-1})).
\end{equation}
For each $x$ such that $xgx^{-1} \in \para_\mu$, consider the corresponding summand in \eqref{expandedB_eqn_summation}.
We repeat a similar argument to the one presented toward the end of the proof of Lemma~\ref{parabolic_counting_lemma}.
Note that the characteristic polynomials of
\(
\{\pi_{j}^\mu (xgx^{-1})
: j \in \N, \, \mu_j = 1\}\)
have degree 1 and hence are a permutation of
\(
\{h_j : j \in \N, \, \mu_j = 1\}.
\)
Likewise, the characteristic polynomials of
\(
\{\pi_{j}^\mu (xgx^{-1}) : j \in \N, \, \mu_j = 2\}\)
have degree 2 and hence are a permutation of
\(
\{h_j : j \in \N, \, \mu_j = 2\},
\)
as there are no degree-1 factors remaining.
Continuing, we see that, for each \(s \in \N\), the degree-\(s\) characteristic polynomials of the diagonal blocks of \(xgx^{-1}\) are a permutation of the degree-\(s\) irreducible factors of the characteristic polynomial of \(g\).
Moreover, the value of each \(P_s^{\ell_f \cdot \alpha_d(s)}\) in \eqref{expandedB_eqn_summation} depends only on the characteristic polynomial of the argument.
This implies that the product of the Primary-support characters in \eqref{expandedB_eqn_summation} is constant over the sum.
Therefore,
\begin{equation}
\label{pre_near_the_end}
\chi^{f \mapsto \lambda} (g) = (-1)^{\tfrac{n}{d} (d-1)} \cdot \frac{\chi^\lambda_{\tilde{\mu}}}{z_{\tilde{\mu}}} \cdot \frac{\# \{x \in \GL_n \F_q : xyx^{-1} \in \para_\mu \} }{\# \para_{\mu}} \cdot \prod_{i = 1}^{\ell(\mu)} P_{\mu_i}^{\ell_f \cdot \alpha_d (\mu_i)} (\pi^\mu_i (g)).
\end{equation}
By Lemma \ref{parabolic_counting_lemma}, \eqref{pre_near_the_end} reduces to
\begin{equation}
\label{near_the_end}
\chi^{f \mapsto \lambda} (g) = (-1)^{\tfrac{n}{d}(d-1)} \cdot
\chi^\lambda_{\tilde{\mu}}
\cdot
\prod_{i = 1}^{\ell(\mu)}
\frac{1}{\tilde{\mu}_i}
P_{\mu_i}^{\ell_f \cdot \alpha_d (\mu_i)} (\pi^\mu_i(g)).
\end{equation}

Consider the Primary-support character evaluations in \eqref{near_the_end}.
By \eqref{who_are_my_roots} and definition \eqref{P_chars},
\begin{align}
P_{\mu_i}^{\ell_f \cdot \alpha_d (\mu_i)} (\pi^\mu_i (g))
& = \sum_{j = 1}^{\mu_i} \te(\eps_{\mu_i}^{\ell_{h_i}})^{q^j \ell_f \cdot [\tilde{\mu}_i]_{q^d}}
= \sum_{\substack{\beta_i \in \F_{q^{\mu_i}} \\ h_i( \beta_i) = 0}} \te ( \beta_i )^{\ell_f \cdot [\tilde{\mu}_i]_{q^d}}. \label{very_last_step}
\end{align}
Substituting \eqref{very_last_step} into \eqref{near_the_end}
gives the result.
\end{proof}

\begin{example}
Returning to Example~\ref{runningexample}, we compute $\chi^{f_3 \mapsto (1)} (c)$ for $c \in E$.
In the notation of Theorem \ref{main_lemma_in_intro}, we have $d = 3, f = f_3, \lambda = (1), \mu = (3), \tilde{\mu} = (1), h_1 = f_3$, and $\ell_f = 1$.
The roots of $h_1$ are $\eps_3, \eps_3^2$, and $\eps_3^4$.
By Theorem \ref{main_lemma_in_intro},
\begin{align*}
\chi^{f_3 \mapsto (1)} (c)
& =
(-1)^{\frac{3}{3}(3-1)}
\cdot
\chi^{(1)}_{(1)}
\cdot
\sum_{h_1(\beta) = 0} \te(\beta) \\
& = \te(\eps_3) + \te(\eps_3^2) + \te(\eps_3^4)
= \zeta_7 + \zeta_7^2 + \zeta_7^4.
\end{align*}
We see this exact value in Table~\ref{GL3F2_char_table} above.
\end{example}






